\documentclass[11pt]{article}

\usepackage[utf8]{inputenc}
\usepackage[T1]{fontenc}
\usepackage{lmodern}
\usepackage{geometry}
\usepackage{amsmath,amssymb,amsfonts,amsthm,mathtools}
\usepackage{bm}
\usepackage{enumitem}
\usepackage{microtype}
\usepackage{hyperref}
\usepackage{titlesec}
\usepackage{booktabs}
\usepackage{array}
\usepackage{setspace}
\usepackage{mathrsfs}
\usepackage{physics}

\geometry{margin=1in}
\hypersetup{colorlinks=true, linkcolor=black, urlcolor=blue, citecolor=black}
\setlist[enumerate]{topsep=0.4em,itemsep=0.35em}
\setlist[itemize]{topsep=0.3em,itemsep=0.25em}
\titleformat{\section}{\large\bfseries}{\thesection.}{0.5em}{}
\titleformat{\subsection}{\normalsize\bfseries}{\thesubsection.}{0.5em}{}
\setstretch{1.06}

\newcommand{\Aether}{\AE{}ther}
\newcommand{\Aflow}{\AE{}ther-flow}
\newcommand{\TheoryName}{\Aflow{} Interpretation of Relativity}
\newcommand{\TheoryProgram}{\Aether{} / \Aflow{} framework}
\newcommand{\ObserverMetric}{g^{\mathrm{br}}}
\newcommand{\CandidateMetric}{g^{\mathrm{cand}}}
\newcommand{\CompletedMetric}{g^{\sharp}}
\newcommand{\PresentedMetric}{\widetilde g}
\newcommand{\RedefMetric}{\widehat g}
\newcommand{\BaseShift}{\Sigma^{\mathrm{IER}}}
\newcommand{\ResponseShift}{\Sigma^{\mathrm{resp}}}
\newcommand{\CompShift}{\Sigma^{\mathrm{cmp}}}
\newcommand{\ResidualCore}{\mathcal V_{\mu\nu}^{\mathrm{res}}}
\newcommand{\ResidualDec}{\mathcal D_{\mu\nu}^{\mathrm{res}}}
\newcommand{\MixQuad}{\mathcal M_{\mu\nu}^{\mathrm{mix}}}
\newcommand{\RespQuad}{\mathcal Q_{\mu\nu}^{\mathrm{resp}}}
\newcommand{\DeltaTCand}{\Delta T_{\mu\nu}^{\mathrm{cand}}}
\newcommand{\ReducedEin}{\mathcal L_{\ObserverMetric}}
\newcommand{\GaugeBook}{\mathcal G_{\ObserverMetric}}
\newcommand{\CompMix}{\mathcal M_{\mu\nu}^{\mathrm{cmp}}}
\newcommand{\CompQuad}{\mathcal Q_{\mu\nu}^{\mathrm{cmp}}}
\newcommand{\DeltaTComp}{\Delta T_{\mu\nu}^{\mathrm{cmp}}}
\newcommand{\ObsBurden}{\mathfrak B_{\mu\nu}^{\mathrm{obs}}}
\newcommand{\ObserverClass}[1]{\left[#1\right]_{\mathrm{obs}}}

\newtheorem{definition}{Definition}
\newtheorem{proposition}{Proposition}
\newtheorem{remark}{Remark}
\newtheorem{corollary}{Corollary}
\newtheorem{theorem}{Theorem}

\title{\textbf{The \TheoryName{}}\\[0.4em]
\large Substrate-Response / Self-Response Charge-Polarization Candidate\\
\large Exact-Preservation Observer-Equation Fixed-Package Strict-Closure Necessity / Maximality Theorem}
\author{}
\date{}

\begin{document}

\maketitle

\begin{abstract}
The active one-metric exact-preservation observer-equation line has now been carried through both ordered closure-upgrade tests on the currently recorded fixed package. Direct exact absorbability of
\[
\ObsBurden
=
\ResidualDec+\CompMix+\CompQuad-8\pi G_N\,\DeltaTComp
\]
is not yet established there, and the explicit same-package field-redefinition route is also not yet established there. The next honest primary manuscript is therefore no longer another status note. It must prove why the current fixed package cannot be upgraded further in the required strict observer-side-closure sense.

The present manuscript does that in a disciplined package-level form. It proves a necessity theorem first: any admissible same-package strict observer-side closure upgrade on the fixed one-metric package must supply one of exactly two ingredients. Either the completed metric \(\CompletedMetric\) itself closes in the exact observer quotient, which is precisely direct absorbability of \(\ObsBurden\), or the closure theorem proceeds through a nontrivial rewritten one-metric variable, in which case it is precisely an explicit same-package field-redefinition theorem. There is no third same-package observer-side upgrade route on the fixed package.

That exhaustion theorem converts the earlier two test verdicts into a current maximality result. Since neither direct absorbability nor an explicit same-package field-redefinition theorem is presently established on the fixed package, no further strict observer-side closure upgrade is available from that package alone at current recorded scope. Exact observer-side effective equivalence is therefore maximal on the present fixed package with respect to strict same-package one-metric observer-side closure. This is a fixed-package necessity / maximality theorem, not a stronger first-principles derivation claim and not a universal impossibility theorem against every future sharpened package.
\end{abstract}

\section{Introduction}

The observer-equation sequence on the active exact-preservation branch has now reached a stricter stage than before. The completed observer equation was isolated, the remaining displayed burden was classified, the displayed sectors were shown not to survive as additional independent observer-side physics, and then the two ordered closure-upgrade tests were executed on the same fixed one-metric package. The first test did not establish direct exact absorbability of \(\ObsBurden\). The second test did not establish an explicit same-package field-redefinition theorem.

That sequence changes what the next honest manuscript must do. Another note that merely repeats the same fixed package, or deepens origin while leaving the same observer equation untouched, does not yet answer the live bottleneck. The bottleneck is now sharper:
\begin{quote}
why, on the currently recorded fixed one-metric package itself, is there no further strict observer-side closure upgrade available unless one of the two tested upgrade ingredients is actually supplied?
\end{quote}
The present note answers exactly that question and nothing broader. Its positive content is a fixed-package necessity / maximality theorem. It does not claim a completed first-principles derivation of GR from substrate microphysics, and it does not claim a universal impossibility theorem against every future sharpened or enlarged package.

\paragraph{Framework and claim boundary.}
\TheoryProgram{} denotes the broader \Aether{} / \Aflow{} framework built on the claims that the \Aether{} is the underlying four-dimensional substrate of reality, the \Aflow{} is its intrinsic ordered motion, observed three-dimensional space is the local experiential slice of that deeper substrate, S-time is the experienced order of change arising from matter, light, and the \Aflow{}, and observed expansion is the three-dimensional appearance of deeper four-dimensional ordered motion. Within that framework, \TheoryName{} denotes the exact-closure theory in which the effective gravitational dynamics are adopted to be exactly Einsteinian with universal matter coupling. In the active sequence, \emph{adoption} means use of that established relativistic dynamics without claiming substrate derivation, while \emph{derivation} is reserved for a first-principles recovery from explicit substrate variables.

The present note stays strictly inside that boundary. It does not promote the active derivational line beyond what the current fixed package proves. Its claim is narrower and package specific: the present fixed one-metric observer-side result is already maximal at current scope unless one of the two missing strict-closure upgrade ingredients is genuinely added.

\section{Fixed Package and Strict Target}

\begin{definition}[Currently recorded fixed one-metric exact-preservation package]
\label{def:fixed}
The present note keeps the following observer-level data fixed.
\begin{enumerate}
    \item The candidate and completed operative metrics are
    \begin{equation}
    \CandidateMetric
    =
    \ObserverMetric+\BaseShift+\ResponseShift,
    \qquad
    \CompletedMetric
    =
    \ObserverMetric+\BaseShift+\ResponseShift+\CompShift.
    \label{eq:metrics}
    \end{equation}
    \item The fixed completion solve is
    \begin{equation}
    \ReducedEin(\CompShift)=-\ResidualCore.
    \label{eq:solve}
    \end{equation}
    \item The completed observer equation is
    \begin{equation}
    G_{\mu\nu}[\CompletedMetric]
    =
    8\pi G_N\,T_{\mu\nu}^{\mathrm{matter}}[\CompletedMetric,\psi]
    +\GaugeBook(\CompShift)_{\mu\nu}
    +\ObsBurden,
    \label{eq:completed}
    \end{equation}
    with
    \begin{equation}
    \ObsBurden
    \equiv
    \ResidualDec+\CompMix+\CompQuad-8\pi G_N\,\DeltaTComp.
    \label{eq:burden}
    \end{equation}
    \item The ordinary matter route is still exactly the same one-metric route
    \begin{equation}
    \DeltaTComp
    \equiv
    T_{\mu\nu}^{\mathrm{matter}}[\CompletedMetric,\psi]
    -
    T_{\mu\nu}^{\mathrm{matter}}[\CandidateMetric,\psi].
    \label{eq:deltatcomp}
    \end{equation}
    \item The inherited candidate-branch decomposition
    \begin{equation}
    \ResidualDec
    =
    \MixQuad+\RespQuad-8\pi G_N\,\DeltaTCand
    \label{eq:resdec}
    \end{equation}
    is kept fixed.
    \item Gauge bookkeeping is already extracted explicitly in \eqref{eq:completed}, and the constraint-exact elimination sectors have already been removed before \(\ObsBurden\) is defined.
    \item The strongest previously established exact observer-side theorem on this branch is exact observer-side effective equivalence: the displayed sectors in \(\ObsBurden\) are exactly non-independent, so the completed branch is exactly observer-side effectively equivalent to Einstein--Hilbert plus ordinary matter through the same one operative metric \(\CompletedMetric\).
    \item The direct exact absorbability test has already been executed on this same fixed package, and direct absorbability of \(\ObsBurden\) is not yet established there.
    \item The explicit same-package field-redefinition test has also been executed on this same fixed package, and an explicit same-package field-redefinition upgrade is not yet established there.
\end{enumerate}
\end{definition}

\begin{definition}[Exact observer quotient]
\label{def:quot}
For the fixed package of Definition \ref{def:fixed}, the \emph{exact observer quotient} is the quotient of observer-side tensors by the already-extracted null classes:
\begin{enumerate}
    \item gauge-exact bookkeeping tensors;
    \item constraint-exact elimination tensors.
\end{enumerate}
The observer-quotient class of a tensor \(X_{\mu\nu}\) is denoted \(\ObserverClass{X_{\mu\nu}}\).
\end{definition}

\begin{remark}
Definition \ref{def:quot} is not enlarged by exact non-independence, benchmark decoupling, or any unspecified new rewriting equivalence. The null classes remain exactly the already-extracted gauge-exact and constraint-exact sectors.
\end{remark}

\begin{definition}[Strict observer-side closure on the fixed package]
\label{def:strict}
The fixed package is said to satisfy \emph{strict observer-side closure} if there exists a one operative metric variable \(\PresentedMetric_{\mu\nu}\), determined by the fixed package and carrying the same single operative matter route, such that
\begin{equation}
\ObserverClass{G_{\mu\nu}[\PresentedMetric]-8\pi G_N\,T_{\mu\nu}^{\mathrm{matter}}[\PresentedMetric,\psi]}=0,
\label{eq:strict}
\end{equation}
with no second operative metric, no second matter law, and no enlargement of the null classes of Definition \ref{def:quot}.
\end{definition}

\begin{definition}[Admissible fixed-package strict-closure upgrade]
\label{def:admissible}
An \emph{admissible fixed-package strict-closure upgrade} is a theorem that establishes Definition \ref{def:strict} while keeping the package of Definition \ref{def:fixed} unchanged. In particular, such a theorem:
\begin{enumerate}
    \item uses only the already fixed one-metric branch data;
    \item introduces no second operative metric and no second matter law;
    \item enlarges none of the null classes in Definition \ref{def:quot};
    \item if it employs a presenting metric variable different from \(\CompletedMetric\), still claims strict closure for the same fixed one-metric package rather than for a new enlarged package.
\end{enumerate}
\end{definition}

\section{Necessity of the Two Upgrade Ingredients}

\begin{proposition}[Closure through \texorpdfstring{\(\CompletedMetric\)}{g-sharp} is exactly direct absorbability]
\label{prop:directequiv}
On the fixed package, the strict observer-side closure statement \eqref{eq:strict} with \(\PresentedMetric=\CompletedMetric\) is equivalent to
\begin{equation}
\ObserverClass{\ObsBurden}=0.
\label{eq:directzero}
\end{equation}
Therefore strict closure through the original completed metric is exactly the direct absorbability route.
\end{proposition}

\begin{proof}
Set \(\PresentedMetric=\CompletedMetric\) in \eqref{eq:strict}. By the completed observer equation \eqref{eq:completed},
\[
G_{\mu\nu}[\CompletedMetric]-8\pi G_N\,T_{\mu\nu}^{\mathrm{matter}}[\CompletedMetric,\psi]
=
\GaugeBook(\CompShift)_{\mu\nu}
+\ObsBurden.
\]
Taking the observer quotient of Definition \ref{def:quot} removes the already-extracted gauge bookkeeping term \(\GaugeBook(\CompShift)\). Hence
\[
\ObserverClass{G_{\mu\nu}[\CompletedMetric]-8\pi G_N\,T_{\mu\nu}^{\mathrm{matter}}[\CompletedMetric,\psi]}
=
\ObserverClass{\ObsBurden}.
\]
Therefore \eqref{eq:strict} with \(\PresentedMetric=\CompletedMetric\) holds if and only if \eqref{eq:directzero} holds. This is exactly direct absorbability of the burden in the exact observer quotient.
\end{proof}

\begin{definition}[Explicit same-package field-redefinition theorem]
\label{def:redef}
An \emph{explicit same-package field-redefinition theorem} on the fixed package means an explicit theorem exhibiting a rewritten metric variable \(\RedefMetric_{\mu\nu}\) such that:
\begin{enumerate}
    \item \(\RedefMetric_{\mu\nu}\) is determined by the already fixed package data;
    \item the completed observer equation is rewritten exactly in terms of \(\RedefMetric_{\mu\nu}\) with no second operative metric and no second matter law;
    \item the rewriting is invertible at the observer-side package level, so the theorem remains a theorem about the same fixed package rather than a replacement package;
    \item the rewritten observer equation satisfies
    \[
    \ObserverClass{G_{\mu\nu}[\RedefMetric]-8\pi G_N\,T_{\mu\nu}^{\mathrm{matter}}[\RedefMetric,\psi]}=0
    \]
    modulo only the already-extracted null classes.
\end{enumerate}
\end{definition}

\begin{proposition}[Any non-direct admissible upgrade is a same-package field-redefinition theorem]
\label{prop:redef}
Let an admissible fixed-package strict-closure upgrade establish \eqref{eq:strict} through a presenting metric \(\PresentedMetric\neq\CompletedMetric\). Then that theorem necessarily supplies an explicit same-package field-redefinition theorem in the sense of Definition \ref{def:redef}.
\end{proposition}

\begin{proof}
By Definition \ref{def:admissible}, the theorem still concerns the same fixed one-metric package and introduces no second operative metric and no second matter law. Since \(\PresentedMetric\neq\CompletedMetric\), the presenting metric is a rewritten variable built from the same fixed package data. Because the theorem claims strict closure for the \emph{same} package rather than a new enlarged package, the rewritten observer equation must be exactly equivalent to the original completed observer equation at the observer-side package level; otherwise the theorem would have changed the package instead of rewriting it. That exact equivalence is precisely the needed observer-side invertibility. Therefore the admissible non-direct route is exactly an explicit same-package field-redefinition theorem.
\end{proof}

\begin{theorem}[Necessity / exhaustion theorem for fixed-package strict closure]
\label{thm:necessity}
Every admissible fixed-package strict-closure upgrade must supply at least one of the following two ingredients:
\begin{enumerate}
    \item the direct absorbability identity \(\ObserverClass{\ObsBurden}=0\);
    \item an explicit same-package field-redefinition theorem in the sense of Definition \ref{def:redef}.
\end{enumerate}
Equivalently, there is no third same-package observer-side route from exact observer-side effective equivalence to strict observer-side closure on the fixed one-metric package.
\end{theorem}

\begin{proof}
Let an admissible fixed-package strict-closure upgrade be given. By Definition \ref{def:strict}, it establishes \eqref{eq:strict} for some one operative metric variable \(\PresentedMetric\) on the same fixed package.

If \(\PresentedMetric=\CompletedMetric\), Proposition \ref{prop:directequiv} shows that the theorem is exactly the direct absorbability route and therefore supplies item (1).

If \(\PresentedMetric\neq\CompletedMetric\), Proposition \ref{prop:redef} shows that the theorem is exactly an explicit same-package field-redefinition theorem and therefore supplies item (2).

These two cases exhaust all admissible fixed-package strict-closure upgrades. Hence no third same-package route exists.
\end{proof}

\begin{remark}
Theorem \ref{thm:necessity} is the necessity part of the present manuscript. It does not yet use the earlier two negative test verdicts. It first proves the structural point that any successful strict same-package observer-side closure theorem on the fixed package must contain one of exactly the two missing upgrade ingredients already isolated in the preceding sequence.
\end{remark}

\section{Current Fixed-Package Maximality / No-Go Result}

\begin{theorem}[Current fixed-package strict-closure maximality theorem]
\label{thm:maximality}
On the currently recorded fixed one-metric exact-preservation package, no admissible fixed-package strict-closure upgrade is presently available. Equivalently, the exact observer-side effective-equivalence theorem is maximal on the current fixed package with respect to strict same-package one-metric observer-side closure.
\end{theorem}

\begin{proof}
By Definition \ref{def:fixed}(8), the direct exact absorbability test has already been executed on the current fixed package and does not establish \(\ObserverClass{\ObsBurden}=0\) there. By Definition \ref{def:fixed}(9), the explicit same-package field-redefinition test has also been executed and does not establish an explicit same-package field-redefinition theorem there. But Theorem \ref{thm:necessity} proves that every admissible fixed-package strict-closure upgrade must supply at least one of exactly those two ingredients. Therefore no admissible fixed-package strict-closure upgrade is presently available on the currently recorded package.

Definition \ref{def:fixed}(7) already records the strongest positive observer-side theorem presently proved on that package, namely exact observer-side effective equivalence. Since no further admissible strict-closure upgrade is presently available, that effective-equivalence theorem is maximal on the current fixed package with respect to strict same-package one-metric observer-side closure.
\end{proof}

\begin{corollary}[Current no-third-route no-go on the fixed package]
\label{cor:nothird}
At current recorded scope, a manuscript that merely restates, reorganizes, or deepens the same fixed one-metric observer package without supplying either a stronger direct absorbability identity or an explicit same-package field-redefinition theorem cannot upgrade the branch to strict observer-side closure.
\end{corollary}

\begin{proof}
Theorem \ref{thm:necessity} removes the possibility of a third same-package route. Theorem \ref{thm:maximality} then records that neither of the only two admissible upgrade ingredients is presently established on the current fixed package. Therefore a manuscript that preserves the same fixed package but supplies neither missing ingredient cannot produce the required strict observer-side closure upgrade.
\end{proof}

\begin{corollary}[Primary-status rule for deeper-origin continuation after maximality]
\label{cor:depth}
After Theorem \ref{thm:maximality}, a deeper-origin continuation is primary on this observer-side bottleneck only if it directly supplies one of the two missing strict-closure ingredients: either a stronger direct absorbability identity or an explicit same-package field-redefinition theorem. A deeper-origin continuation that merely preserves the same completed observer equation is supporting context rather than the immediate next primary move.
\end{corollary}

\begin{proof}
Corollary \ref{cor:nothird} shows that preserving the same fixed observer package without one of the two missing ingredients cannot upgrade the branch to strict observer-side closure. Therefore deeper-origin continuation changes the live primary burden only if it directly furnishes one of those two ingredients.
\end{proof}

\section{What This Note Establishes and What It Does Not}

The present note establishes the following.
\begin{enumerate}
    \item It proves a necessity theorem for the current fixed one-metric package: any admissible strict observer-side closure upgrade on that package must supply either direct absorbability of \(\ObsBurden\) or an explicit same-package field-redefinition theorem.
    \item It proves that there is no third same-package observer-side route from exact observer-side effective equivalence to strict observer-side closure on the fixed package.
    \item It converts the already recorded failure of both ordered upgrade tests into a current fixed-package maximality result: exact observer-side effective equivalence is maximal on the currently recorded package with respect to strict same-package one-metric observer-side closure.
\end{enumerate}

The present note does \emph{not} establish the following.
\begin{enumerate}
    \item It does not claim a completed first-principles derivation of GR from substrate microphysics.
    \item It does not claim a universal impossibility theorem against every future sharpened or enlarged package.
    \item It does not forbid deeper-origin work that directly supplies one of the two missing upgrade ingredients.
    \item It does not enlarge the observer quotient beyond the already-extracted gauge-exact and constraint-exact sectors.
\end{enumerate}

\begin{remark}
Theorem \ref{thm:maximality} is therefore a fixed-package maximality / no-go statement, not a universal no-go against every future path in the active derivational program. Its force is narrower and exactly the force now needed: the current fixed package has no further honest strict observer-side closure upgrade available unless one of the two missing upgrade ingredients is genuinely added.
\end{remark}

\section{Conclusion}

The two ordered closure-upgrade tests on the current one-metric observer-equation bottleneck have now both been executed, and neither is presently established on the fixed package. The present manuscript proves why that fact now has theorem-level force rather than merely project-management force.

First, it proves necessity. Any admissible strict observer-side closure theorem on the same fixed package must proceed in one of exactly two ways: either the original completed metric \(\CompletedMetric\) closes directly in the observer quotient, which is precisely direct absorbability of \(\ObsBurden\), or the theorem proceeds through a rewritten one-metric variable, which is precisely an explicit same-package field-redefinition theorem. There is no third same-package observer-side route.

Second, it proves maximality at current recorded scope. Because the direct absorbability test and the same-package field-redefinition test have both already been carried out on the fixed package and neither ingredient is presently established there, the current package cannot be upgraded further in the required strict observer-side-closure sense by its current recorded data alone. The strongest exact observer-side result on that package therefore remains exact observer-side effective equivalence.

This is the positive result, the scope limit, and the next honest burden. The positive result is a fixed-package necessity / maximality theorem. The scope limit is that this is not a universal impossibility theorem and not a derivation claim. The next honest burden is equally sharp: any deeper or sharper continuation becomes primary on this bottleneck only if it directly supplies a stronger absorbability identity or an explicit same-package field-redefinition theorem.

\end{document}
